\section{问题二：固定电价下的因果滚动购电与储能调度}
\label{sec:q2}

\subsection{任务与信息条件}
问题二的分时电价按日内时刻重复。负载与光伏的未来实际值在决策时刻不可见，正常购电合同须在每日0:00一次性提交144个十分钟区间，日内不再修改；储能则每十分钟滚动决策，只执行当前动作。若冻结合同、实际光伏利用和储能放电仍不能覆盖当期负载及已承诺充电，则以该时刻正常电价的5倍紧急购电，未完全使用的正常合同仍按合同电量计费。

2025年1月作为历史初始化期，储电量保持6\,000 kWh；自2月1日起连续回放至12月31日。实际负载、光伏仅在区间结束后进入历史，因此预测、日前合同、实际执行与实际账单在记录中严格区分。

\subsection{因果预测与滚动优化}
\subsubsection{负载与光伏预测}
对日$d$、区间$k$，负载先取过去7、14、21、28日同刻值的等权周周期基线，再叠加最近四个已结束日的同刻预测残差均值：
\begin{equation}
 \widehat\ell_{d,k}=
 \frac14\sum_{j=1}^{4}\ell_{d-7j,k}
 +0.8\cdot\frac14\sum_{i=1}^{4}
 \left(\ell_{d-i,k}-\widetilde\ell_{d-i,k}\right),\qquad
 \widehat\ell_{d,k}\leftarrow\max\{0,\widehat\ell_{d,k}\}.
 \label{eq:q2-load-forecast}
\end{equation}
其中$\widetilde\ell$为当日保存的原预测，残差修正不读取未来实际值。光伏采用过去1--5日同刻实际出力的递减加权组合：
\begin{equation}
 \begin{aligned}
 \widehat v_{d,k}&=\max\left\{0,\sum_{i=1}^{5}\alpha_i v_{d-i,k}\right\},\\
 (\alpha_1,\ldots,\alpha_5)&=(0.35,0.25,0.20,0.12,0.08).
 \end{aligned}
 \label{eq:q2-pv-forecast}
\end{equation}
保存的全年预测记录显示，负载、光伏RMSE分别约为182 kW和305 kW。该误差仅描述可见信息下的预测偏差；最终费用仍以实际回放账单计算。

\subsubsection{日前合同与十分钟储能动作}
设$g_{t+h}$、$e^{\mathrm p}_{t+h}$、$c_{t+h}$、$d_{t+h}$、$u_{t+h}$分别为24小时滚动窗口内的正常合同、规划紧急购电、母线侧充电、放电和光伏利用量。以固定电价$p_{t+h}$最小化新增合同费、预测缺口费与窗口末态价值：
\begin{equation}
 \min J_t=\sum_{h\in\mathcal H_t}
 \left(p_{t+h}g_{t+h}+5p_{t+h}e_{t+h}^{\mathrm p}\right)
 -\nu_tE_{\mathrm{end}}.
 \label{eq:q2-objective}
\end{equation}
窗口达到年度终点时，移除末态价值并直接约束年末储电量为6\,000 kWh。规划供需、储能状态与互斥约束为
\begin{align}
 g_{t+h}+e_{t+h}^{\mathrm p}+d_{t+h}+u_{t+h}
 &\ge\widehat\ell_{t+h}+c_{t+h},\\
 0&\le u_{t+h}\le\widehat v_{t+h},\\
 E_{t+h+1}&=E_{t+h}+0.9c_{t+h}-d_{t+h}/0.9,\\
 1200&\le E_{t+h}\le10800,\\
 0\le c_{t+h}&\le Mz_{t+h},\qquad
 0\le d_{t+h}\le M(1-z_{t+h}),\\
 M&=\frac{5000}{6},\qquad z_{t+h}\in\{0,1\}.
 \label{eq:q2-model}
\end{align}
每日0:00将该日$g_{d,k}$写入合同账本并冻结；后续优化只更新储能后继动作，不改写当日合同。模型采用SciPy的\texttt{milp}接口调用HiGHS求解\cite{scipyMilp,highs}，完整流程见图\ref{fig:q2-workflow}。

\begin{figure}[H]
  \centering
  \PaperGraphic[width=0.98\linewidth]{figures/q2_workflow.pdf}
  \caption{问题二的因果预测、合同冻结与实际回放流程}
  \label{fig:q2-workflow}
\end{figure}

\subsection{实际执行与费用结算}
每次滚动仅执行当前求得的充、放电动作，实际储电量按式\eqref{eq:q2-model}连续传递。设实际负载、冻结合同、实际光伏利用和已执行充放电为$\ell_t,g_t,u_t,c_t,d_t$，实际紧急购电定义为
\begin{equation}
 e_t=\max\{\ell_t+c_t-g_t-u_t-d_t,0\}.
 \label{eq:q2-replay-emergency}
\end{equation}
式中$c_t$为已承诺并执行的充电动作，故紧急购电可以弥补执行该动作后仍存在的即时缺口；它不同于预测优化中的$e_t^{\mathrm p}$。最终账单由
\begin{align}
 F_{\mathrm{plan}}&=\sum_tp_tg_t,&
 F_{\mathrm{emergency}}&=\sum_t5p_te_t,&
 F_{\mathrm{total}}&=F_{\mathrm{plan}}+F_{\mathrm{emergency}}
 \label{eq:q2-cost}
\end{align}
计算，不累计滚动预测目标。

\subsection{回放结果与分析}
选定的非合成运行覆盖2025年2月1日至12月31日共334日、48\,096个区间。全年正常合同费为1\,350.0196万元，实际紧急购电费为168.5776万元，总费用为1\,518.5972万元；实际紧急购电量为47.1974万kWh，见表\ref{tab:q2-cost}。

\begin{table}[H]
\centering
\small
\caption{问题二连续回放的实际费用}\label{tab:q2-cost}
\begin{tabular}{lrrr}
\toprule
回放范围 & 正常合同费（万元） & 紧急购电费（万元） & 总费用（万元）\\
\midrule
2025-02-01至2025-12-31 & 1350.0196 & 168.5776 & 1518.5972\\
\bottomrule
\end{tabular}
\end{table}

图\ref{fig:q2-costs}(a)表明紧急费用并非集中于单月，而是由冻结合同与实际供需偏差共同形成。图\ref{fig:q2-costs}(b)直接比较实际账单与任务目标：总费用较1\,500万元高18.5972万元（1.2398\%），紧急购电费较100万元高68.5776万元。因此，本次结果完整反映当前机制的真实回放表现，但未同时达到两项费用目标，不能以计划费用或预测目标替代实际账单。

\begin{figure}[H]
  \centering
  \PaperGraphic[width=0.98\linewidth]{figures/q2_costs.pdf}
  \caption{问题二的月度实际费用构成及年度目标比较}
  \label{fig:q2-costs}
\end{figure}

\subsubsection{紧急购电的分布与形成机制}
紧急购电并非由少数异常日主导。334个回放日中有333日发生紧急购电，而紧急购电量最高的10日合计仅占全年紧急电量的13.91\%。从日内分布看，10:00--16:00贡献全年紧急电量的59.35\%，其中14:00--16:00两小时即占24.40\%。因此，168.5776万元紧急购电费反映的是重复出现的结构性缺口，而非个别日期的偶然失策。

将每日0:00保存的日前预测与48\,096个实际区间逐一配对后发现：在发生紧急购电的7\,961个区间内，光伏预测平均高于实际206.91 kW，负载预测平均低于实际160.32 kW，二者共同使预测净负荷比实际净负荷低367.23 kW；在10:00--16:00的紧急区间内，光伏高估进一步扩大到489.07 kW，而负载低估为74.58 kW，对应净负荷低估563.65 kW。若按相同日内时刻比较不同日期，紧急区间的实际光伏仍平均低138.01 kW，说明上述结论并非仅由紧急购电集中在白天所造成。全年全部区间的光伏预测平均偏差仅为$+1.69$ kW，故这里所称的“系统性”是指不利误差尾部在几乎全年反复触发，而不是断言预测在所有时段均持续偏高。

该现象可由“日前冻结---净负荷低估---高价补缺”解释。首先，正常合同在0:00按当日预测一次性冻结，日内不能补签；其次，光伏受天气驱动，当前1--5日同刻外推在午间仍会出现较大的单侧误差，同时负载也可能被低估；再次，统一5\%负荷裕度只是固定比例缓冲，不能完全覆盖上述联合误差尾部。当冻结合同和受功率、容量约束的储能均无法填补实际缺口时，式\eqref{eq:q2-replay-emergency}使剩余误差全部进入紧急购电。因此，本方案不是主观“冒险少买”，更准确的表述是：模型以点预测为基础、只配置固定安全裕度，并在期望费用目标下承担预测误差尾部；100万元是结果评价指标，而不是模型中的可行性硬约束。

紧急购电量47.1974万kWh仅相当于实际利用光伏电量1680.7935万kWh的2.81\%，但紧急平均单价约为3.57元/kWh，使其费用占总成本的11.10\%。若要进一步压低紧急购电，需在日前增加不可退的正常合同，或引入随误差分布变化的分时安全裕度和情景采购；这会同时增加正常合同费以及合同余电风险。只有当节省的5倍紧急费用大于新增合同费用时，总成本才会下降，故不能仅以紧急费用下降判断策略优劣，最优裕度仍需通过独立敏感性回放确定。

此外，账本给出的226.0955万kWh“未利用光伏”与47.1974万kWh紧急购电在年度汇总中并存并不矛盾。逐区间复核显示，二者没有在同一十分钟区间同时出现：光伏残余主要发生在供给富余时段，紧急购电则发生在实际净负荷高于冻结合同和可用储能的时段；受储能容量、功率和效率限制，前一时段的光伏不能无损搬移至后一时段。并且本文采用\texttt{grid\_and\_discharge\_first\_residual\_pv}归因口径，先以合同电量和放电冲抵母线需求，再将剩余需求归因于光伏利用，故226.0955万kWh是该结算顺序下的光伏残余，不应脱离此口径简单解释为物理上主动弃光。

上述诊断表明，紧急购电超出100万元目标的直接原因，是日前冻结机制将净负荷预测的单侧误差转化为唯一的高价补缺通道；其中午间光伏高估最为突出，固定5\%负荷裕度尚不足以完全吸收该误差尾部。

图\ref{fig:q2-day}选取全年紧急购电量最大的2025年6月1日。图中正常购电为日初冻结合同，充放电为逐十分钟实际执行量；日内多次出现已执行充电和紧急购电并存，这正是式\eqref{eq:q2-replay-emergency}将已承诺充电计入即时需求的结果，而逐区间充放电仍保持互斥。

\begin{figure}[H]
  \centering
  \PaperGraphic[width=0.94\linewidth]{figures/q2_2025-06-01.pdf}
  \caption{问题二在2025年6月1日的实际回放轨迹}
  \label{fig:q2-day}
\end{figure}

\subsection{复核与局限}
保存的验证器覆盖全部48\,096个区间，最大约束违反量为0 kWh；按账本重新计算的最大电量残差为$4.55\times10^{-13}$ kWh，光伏分配残差为0。储电量从期初6\,000 kWh连续传递至期末并恢复为6\,000 kWh，逐区间未出现同时充、放电。

\begin{table}[H]
\centering
\small
\caption{问题二实际回放的复核结果}\label{tab:q2-check}
\begin{tabular}{lll}
\toprule
检验项目 & 计算结果 & 判定\\
\midrule
回放区间数 & 48\,096 & 通过\\
最大约束违反量 & 0 kWh & 通过\\
最大账本电量残差 & $4.55\times10^{-13}$ kWh & 通过\\
光伏分配残差 & 0 kWh & 通过\\
期初、期末储电量 & 6\,000、6\,000 kWh & 通过\\
逐区间充放电互斥 & 未出现同时充放电 & 通过\\
\bottomrule
\end{tabular}
\end{table}

模型的限制在于仅用历史实际值构造预测，异常天气或负载变化会直接传导至冻结合同；普通合同又不能日内修正，储能功率、容量不足时只能以高价紧急购电。费用也未计电池寿命损耗。因此，本文报告当前模型与数据下的实际账单，而非对任意未来场景的费用保证。
\FloatBarrier
